\section*{第 \ref{ch:g2:moon-in-the-sky} 章 --- 天空中的月亮}

\begin{solution}{exo:g2:moon-in-the-sky:1}
比地球小得多，比云远得多 ---
比任何飞机飞得都高。我们知道人可以去，
因为宇航员去过：单程要三天。
\end{solution}

\begin{solution}{exo:g2:moon-in-the-sky:2}
环形山 --- 碗状伤痕，是太空中的岩石撞上月球留下的。
它们之所以能保留下来，
是因为月球上没有风，也没有雨来把它们磨平。
\end{solution}

\begin{solution}{exo:g2:moon-in-the-sky:3}
光的收集者。月光是太阳光从月球灰色的岩石上反射回来的，
就像白墙能照亮房间，它本身却不是灯。
\end{solution}

\begin{solution}{exo:g2:moon-in-the-sky:4}
总是面向太阳的那一半 --- 朝向太阳的一面明亮，
背面则处于它自己的夜晚。
\end{solution}

\begin{solution}{exo:g2:moon-in-the-sky:5}
能。许多个下午，淡淡的月亮挂在蓝天上。
要让怀疑的朋友信服，最好的证据是：在这样的下午走到户外，
指给他看。
\end{solution}

\begin{solution}{exo:g2:moon-in-the-sky:6}
蛾眉月、半月、近满月、满月 --- 然后再经过近满月、半月、蛾眉月，
回到看不见的“朔”月。
整个过程大约需要一个月。
\end{solution}

\begin{solution}{exo:g2:moon-in-the-sky:7}
大声喊也听不到声音 --- 声音需要空气才能传播。风筝也飞不起来：
没有空气就没有风，风筝也就没有可以借力的东西。
\end{solution}

\begin{solution}{exo:g2:moon-in-the-sky:8}
每天晚上在同一时刻、用同样诚实的画法作画（同一地点也有帮助）：
一份公正的日记每次只改变一件事 ---
夜晚。
\end{solution}

\begin{solution}{exo:g2:moon-in-the-sky:9}
房屋离得近，驶过时萝莎看到它们的样子会改变 ---
房子一掠而过，就不见了。月球非常遥远，整段车程都不会改变她看月亮的样子：
月亮始终在她肩旁，好像在跟着她走。
远处事物“跟随”的现象最为典型。
\end{solution}

\begin{solution}{exo:g2:moon-in-the-sky:10}
是地球本身。阳光从明亮的地球反射到月球的暗面上，
使那里隐隐发亮 --- 这叫地照光：我们的行星
也在回报月球的恩惠。
\end{solution}
